\newpage
\chapter*{\centering Abstract}
\textit{\quad 
This report presents a fairness audit of a high-performing ensemble machine learning model developed for depression prediction, originally designed for a Kaggle competition using synthetic mental health data. While the model achieves high predictive accuracy by combining AdaBoost, CatBoost, XGBoost, Logistic Regression, and LGBMClassifier, our audit reveals critical disparities in subgroup performance. Through detailed analysis of feature correlations, missing data patterns, and regional representation, we identify structural and cultural biases in the dataset, particularly affecting employment status and age subgroups. Stakeholder perspectives (patients, families, professionals, and institutions) are also mapped to understand real-world implications. Fairness evaluations using Fairlearn’s MetricFrame and ThresholdOptimizer expose significant performance gaps: students experience high false positive rates, while working professionals face high false negative rates. Post-processing mitigation with equalized odds constraints reduces these disparities, though at a cost to overall accuracy. LIME is further applied to interpret individual predictions, uncovering group-specific reasoning patterns. This audit highlights the need for fairness-aware evaluation, culturally grounded feature design, and subgroup-specific calibration when deploying automated decision systems in high-stakes domains like mental health.
}

